% Section 2 — Architecture
% Paper: "The Learnable Bridge: Task Fingerprinting and Adapter Composition
%         via Structured Coupling in Low-Rank Adaptation"

\section{Architecture}
\label{sec:architecture}

\subsection{Standard LoRA}
\label{sec:standard-lora}

Low-rank adaptation~\cite{hu2022lora} freezes a pretrained weight matrix
$W_0 \in \mathbb{R}^{d_\text{out} \times d_\text{in}}$ and learns a
low-rank update $\Delta W = BA$, where $B \in \mathbb{R}^{d_\text{out}
\times r}$, $A \in \mathbb{R}^{r \times d_\text{in}}$, and $r \ll
\min(d_\text{in}, d_\text{out})$.
The adapted forward pass is
\begin{equation}
  h = W_0 x + \frac{\alpha}{r}\, B A x,
  \label{eq:lora}
\end{equation}
where $\alpha$ is a fixed scaling factor and $x \in \mathbb{R}^{d_\text{in}}$.
$A$ is initialized via Kaiming uniform~\cite{he2015delving} and $B$ is
initialized to zero, so $\Delta W = 0$ at the start of training.
The rank-$r$ hidden vector $z = Ax \in \mathbb{R}^r$ serves as an
intermediate representation, but LoRA imposes no structure on it:
the $r$ dimensions are an arbitrary basis for a low-rank subspace.
Two dimensions of $z$ interact only insofar as $B$ happens to mix
them in the up-projection---there is no explicit coupling mechanism
within the hidden space itself.

\subsection{The Bridge Modification}
\label{sec:bridge}

We partition the rank-$r$ hidden vector into $C$ channels of width
$r / C$ (requiring $C \mid r$) and insert a learnable \emph{bridge
matrix} $M \in \mathbb{R}^{C \times C}$ that mixes channels before
the up-projection.
Concretely, the adapted forward pass becomes:
\begin{equation}
  h = W_0 x + \frac{\alpha}{r}\, B\,
    \operatorname{vec}\!\bigl(M \cdot
    \operatorname{mat}_{C \times r/C}(Ax)\bigr),
  \label{eq:bridge}
\end{equation}
where $\operatorname{mat}_{C \times r/C}$ reshapes $\mathbb{R}^r
\to \mathbb{R}^{C \times (r/C)}$ and $\operatorname{vec}$ is its
inverse.
Equivalently, in index notation: the hidden vector $z = Ax$ is
reshaped to $Z \in \mathbb{R}^{C \times (r/C)}$, the bridge acts as
$\tilde{Z}_{ik} = \sum_j M_{ij} Z_{jk}$, and $\tilde{Z}$ is
flattened back to $\mathbb{R}^r$ before multiplication by $B$.

In our experiments $r = 24$ and $C = 6$, so each channel carries
$r/C = 4$ rank slices and $M$ is a $6 \times 6$ matrix with 36
learnable parameters.
At a representative target module in Qwen2.5-7B
($d_\text{in} = d_\text{out} = 3{,}584$), the standard LoRA
projection matrices contain $2 \times 24 \times 3{,}584 = 172{,}032$
parameters; the bridge adds 36, a $0.021\%$ overhead.

\paragraph{Identity recovery.}
When $M = I_C$, the reshape--multiply--flatten sequence is the identity
map on $\mathbb{R}^r$.
Equation~\eqref{eq:bridge} reduces exactly to
Equation~\eqref{eq:lora}: the bridge architecture is a strict
generalization of standard LoRA, with LoRA as the $M = I_C$ special
case.
We initialize $M = I_C$ at the start of training so that the bridge
begins inert and any coupling structure in $M$ at convergence is
entirely learned, not inherited from initialization.

\paragraph{Gradient flow.}
The loss gradient with respect to a bridge element $M_{ij}$ is
\begin{equation}
  \frac{\partial \mathcal{L}}{\partial M_{ij}}
  = \sum_{k=1}^{r/C}
    \frac{\partial \mathcal{L}}{\partial \tilde{Z}_{ik}}
    \cdot Z_{jk},
  \label{eq:bridge-grad}
\end{equation}
where $\tilde{Z}_{ik}$ receives its gradient through $B$ and the
downstream loss.
$M_{ij}$ therefore accumulates gradient proportional to the
\emph{co-activation} between channel~$j$ (in the down-projection's
output) and channel~$i$ (in the up-projection's gradient signal).
If training data induces systematic correlations between specific
channel pairs, the corresponding off-diagonal entries of $M$ grow;
if channel activations are independent, $M$ remains near $I_C$.
This selectivity is the mechanism that enables task fingerprinting
(\S\ref{sec:experiments}).

\paragraph{Rotational symmetry.}
For any orthogonal $R \in \mathbb{R}^{r \times r}$, the product $BA$
is invariant under $B \mapsto BR^\top$, $A \mapsto RA$.
The rank dimensions of standard LoRA are therefore rotationally
symmetric---there is no intrinsic ``channel~1'' or ``channel~6.''
The bridge does not break this symmetry: the optimizer may permute or
rotate channels freely, and the bridge will compensate.
What the bridge \emph{does} provide is a readable summary of the
coupling structure that training discovered, regardless of which
rotation the optimizer chose.
Two bridges trained on the same task will not have identical entries,
but they will exhibit the same spectral properties (eigenvalues,
Fiedler connectivity) because these are rotation-invariant.
We exploit this in our analyses by comparing bridges via their
eigenspectra and graph-theoretic connectivity rather than entry-wise
distance.

\subsection{Topological Constraints}
\label{sec:topology}

The bridge $M$ is a dense $C \times C$ matrix by default.
Its off-diagonal entries can optionally be constrained by a binary
mask $\mathcal{T} \in \{0, 1\}^{C \times C}$ derived from a lattice
topology, enforcing $M_{ij} = 0$ wherever $\mathcal{T}_{ij} = 0$.
During training, masked entries receive no gradient and remain at zero.

We define two topological masks for $C = 6$, derived from the
connectivity of three-dimensional lattices:

\paragraph{FCC mask.}
The six channels correspond to the six antipodal face-pairs of the
rhombic dodecahedron---equivalently, the six nearest-neighbor direction
pairs of the face-centered cubic (FCC) lattice.
Two direction pairs \emph{couple} when they share vertices on the
underlying polyhedron.
In the rhombic dodecahedron, co-planar pairs (those lying in the same
coordinate plane, e.g., directions $(1,1,0)$ and $(1,-1,0)$ in the
$xy$-plane) share 4 octahedral vertices; cross-planar pairs share 2.
All pairs share at least 2 vertices, so the FCC mask is the complete
graph on 6 nodes: $\mathcal{T}^\text{FCC}_{ij} = 1$ for all $i, j$.
The mask permits 6 diagonal entries and $\binom{6}{2} = 15$
bidirectional off-diagonal connections, for 36 non-zero entries total.
In practice, the unconstrained bridge is the FCC-masked bridge---the
topology permits full coupling.

\paragraph{Cubic mask.}
The cubic lattice connects each node to 6 neighbors along 3
axis-aligned directions, yielding 3 direction pairs.
To construct a cubic mask on 6 channels, we retain coupling only
within each co-planar pair: channels $(0,1)$, $(2,3)$, and $(4,5)$
may couple, but cross-planar entries are zeroed.
This produces a block-diagonal mask with three $2 \times 2$ blocks:
$\mathcal{T}^\text{cubic}_{ij} = 1$ if $\lfloor i/2 \rfloor =
\lfloor j/2 \rfloor$, and 0 otherwise.
The cubic mask has 6 diagonal entries and 3 bidirectional off-diagonal
connections (one per block), for 12 non-zero entries.

\paragraph{Relationship to MELoRA.}
Setting the mask to $\mathcal{T} = I_C$ (diagonal only, no
off-diagonal coupling) recovers the block-diagonal structure of
MELoRA~\cite{ren2024melora}, where each channel operates independently.
The bridge thus spans a continuum: from fully independent channels
($M = I_C$, MELoRA), through partially coupled (cubic mask), to fully
coupled (FCC mask or unconstrained).
All experiments in this paper use the unconstrained bridge unless
otherwise noted.

\subsection{Diagnostic Quantities}
\label{sec:diagnostics}

The bridge is a small, square, real matrix.
We extract three diagnostic quantities from it, each invariant to the
rotational symmetry of rank dimensions (\S\ref{sec:bridge}):

\paragraph{Fiedler connectivity.}
We interpret $|M_{ij}|$ for $i \neq j$ as the edge weight between
channels $i$ and $j$ in a weighted graph on $C$ nodes, and compute the
algebraic connectivity (second-smallest eigenvalue of the weighted
Laplacian)~\cite{fiedler1973algebraic}.
A Fiedler value of zero indicates a disconnected or identity bridge;
larger values indicate stronger cross-channel coupling.
At initialization ($M = I_C$), Fiedler $= 0$.

\paragraph{Frobenius deviation.}
$\|M - I_C\|_F$ measures how far the bridge has moved from its
identity initialization.
This is a scalar summary of total learned coupling, combining both
diagonal scaling and off-diagonal mixing.

\paragraph{Eigenspectrum.}
The vector of eigenvalues $\lambda_1, \ldots, \lambda_C$ of $M$
(sorted by magnitude) provides a rotation-invariant fingerprint of
the bridge's coupling structure.
We compare eigenspectra between bridges via cosine similarity:
$\operatorname{cos}(\boldsymbol{\lambda}^{(a)},
\boldsymbol{\lambda}^{(b)})$.

\subsection{Figure Description}
\label{sec:arch-figure}

% [Figure 1: Architecture diagram — to be rendered separately]
%
% The figure should show the following, arranged left-to-right as a
% dataflow diagram:
%
% LEFT: Input x (a vertical bar labeled d_in).
%
% BLOCK 1: Down-projection A (r × d_in), producing z ∈ R^r (a
% narrow vertical bar labeled r = 24).
%
% RESHAPE: z is reshaped into a 6 × 4 grid (6 rows = channels,
% 4 columns = rank slices per channel). Each row is color-coded to
% a direction pair. Visual: a small matrix/grid with 6 colored rows.
%
% BLOCK 2: The bridge M (6 × 6), shown as a small square matrix
% overlaid on the channel grid. Arrows or lines connect each
% channel row to the output channel rows, with line thickness
% proportional to |M_{ij}|. The identity initialization is shown
% as a ghosted version (6 diagonal-only connections), with the
% learned state showing off-diagonal connections appearing.
%
% RESHAPE: The mixed 6 × 4 grid is flattened back to R^r.
%
% BLOCK 3: Up-projection B (d_out × r), producing the LoRA
% correction Δh ∈ R^{d_out}.
%
% RIGHT: The correction is added to W_0 x (shown as a separate
% parallel path from x through W_0, merging with Δh via "+").
%
% INSET (bottom-right): Two small 6 × 6 heatmaps side by side:
% (a) M = I_6 (identity, at initialization — diagonal only, all
% off-diagonal white), labeled "Step 0: standard LoRA";
% (b) a learned M (showing off-diagonal structure with varying
% intensity), labeled "Step 10K: learned coupling".
%
% ANNOTATION: A callout on the bridge block reading "36 parameters
% (0.021% of adapter)" with the equation M_{ij}: channel j → channel i.
%
% Caption: "RhombiLoRA architecture. The bridge matrix M couples the
% C = 6 channels of the rank-r hidden representation between the down-
% and up-projections A and B. At initialization M = I_6, recovering
% standard LoRA. Training discovers cross-channel coupling (inset),
% producing a compact diagnostic of adapter behavior."

